\section{SEANCE 2}
L'objectif de cette séance est de réaliser les séries de mesures demandées.
On a le système complet, carte, batterie, capteurs.
On veut ces mesures en particulier

3 mesures sur la batterie, tension au repos, en fonctionnement, 
puis lors de la poussée maximale
donc protocole de test à prévoir, quelle valeur pour la poussée maximale
par rapport à la poussée en fonctionnement normal
Tester les différents états du robot, et à chaque fois, 
mesurer courant et puissance. Les états doivent être définis
couché, balancier, transition, équilibre
mesurer les signaux pwm et la tension du moteur à l'oscilloscope
on veut également les ondulations du bus 5V (à définir, bus etc)
On veut une extrapolation de l'autonomie du robot
donc calcul de l'autonomie

à la fin de cette séance, on veut, un pv de réception, plusieurs schémas,
fonctionnel, structurel, un tableau de mesures qui redonne toutes les valeurs,
pv d'essaies des fonctions d'origine

MESURE 1 & 2
TENSION DE LA batterie
repos, charge (équilibre) poussée maximale (roues bloquées)
COURANT CONSOMMé (ampèremètre en série ou sonde de courant)
calculer puissance appelée par chaque état
MESURE 3 & 4
signaux PWM des moteurs et tension aux bornes d'un moteur
(en avance, équilibre, oscillation volontaire)
MESURE 5
Vérifier la tension du bus 5 V et son ondulation au repos et sous charge
MESURE 6
mesure de la décharge sur une durée d’essai donnée (30
min de fonctionnement modélisé)
CALCULS
Estimer la fréquence de découpage
Estimer l’autonomie énergétique : extrapolation, comparaison à l’usage réel

DOCUMENTS
PV de réception (inventaire, anomalies, références des composants relevées)
Schéma fonctionnel 
(chaînes d’information et d’énergie) et schéma structurel du système
Tableau de mesures puissance : (Vbat , I, P ) pour chaque état, signaux PWM, 
ondulation du bus 5 V, courbe de décharge et estimation d’autonomie commentée
PV d’essais des fonctions d’origine (tableau conforme / non conforme)

On distingue plusieurs états de fonctionnement du robot.
Repos, le robot ne fait aucun mouvement, il est soit sur le dos soit couché.
Transition, le robot passe de couché à équilibre

MESURE 1 & 2 DETAIL
Tout branchement doit être effectué robot éteint (interrupteur batterie arrière).
on veut mesurer la tension de la batterie au repos, en charge, et en poussée maximale.
robot éteint. On place deux câbles sur les pins connectés sur l'entrée de la batterie.
On débranche la batterie, et on vient ralonger la batterie avec deux câbles.
On place une pince ampèremétrique sur le câble relié à la borne + de la batterie.
On téléverse le code donné par le constructeur sur le site de leur produit.
Une fois branché, on le place au sol dans un espace dégagé.
On place un obstacle devant le capteur du robot, son fonctionnement normal est qu'il
va légèrement reculer puis avancer à nouveau, restant alors en équilibre.
On allume le robot avec l'interrupteur situé sur la batterie.
Le robot passe de l'état couché à équilibre, il va chercher à avancer, puis rester 
en équilibre devant un obstacle.
Pour réaliser la mesure au repos, il faut le placer sur le dos, il ne peut donc
pas se mouvoir.
On éteint le robot, on le couche, puis on le rallume. Il va chercher à se relever,
et on réalise la mesure en charge. On éteint le robot.
Pour la mesure à poussée maximale, on vient bloquer les roues.
On éteint le robot une fois la mesure terminée.
On calcule la puissance à partir des résultats, pour chaque état.
On utilise des câbles de longueur adéquate pour que le robot puisse se mouvoir légèrement

MESURE 3 & 4 DETAIL
on place le robot dans les mêmes conditions que pour la première mesure.
On veut mesurer la tension d'alimentation des moteurs, dans tout les cas différents,
ainsi que la tension de commande PWM en sortie de la carte.
Mêmes conditions que la mesure précédente, avec un obsacle devant.
On mesure dans les 3 cas souhaités, repos, équilibre, poussée maximale.

MESURE 5 jsp

MESURE 6 on place le robot en charge (ou lié avec une corde pour ne pas qu'il tombe).
On réalise une mesure de la tension de la batterie, au début de la mesure.
Après 30 minutes, on mesure à nouveau la tension de la batterie.
Avec la différence de potentiel entre les deux instants, on calcule l'autonomie
théorique du robot.

POUR CHAQUE ESSAI
procédure
résultat attendu
résultat constaté

-> reste à faire
trouver noms des broches à tester
leur valeur attendue
